\begin{figure}[htbp]
\centering
\begin{tikzpicture}[>=stealth,scale=1.2]
\draw[->,black!40] (-.3,0)--(3.5,0) node[right] {$x$};
\draw[->,black!40] (0,-.3)--(0,2) node[above] {$y$};
\draw[arkivblaa] (0,0)--(2,0)--(3,1.5)--(1,1.5)--cycle;
\draw[->,arkivblaa] (0,0)--(2,0) node[below] {$u$};
\draw[->,arkivblaa] (0,0)--(1,1.5) node[left] {$v$};
\draw[black!55,dashed] (1,0)--(1,1.5);
\node[right] at (1,.7) {$h$};
\node[anchor=west] at (4,.9) {$\det[u\ v]=u_xv_y-u_yv_x$};
\end{tikzpicture}
\caption{Orienteret areal af et parallelogram: $u=(2,0)$, $v=(1,3/2)$ giver grundlinje 2 og højde $3/2$, altså areal 3. Rækkefølgen $(u,v)$ er positivt orienteret; byttes søjlerne, skifter fortegnet. Akserne er dimensionsløse med samme målestok.}
\label{fig:b1-determinant}
\end{figure}
